\documentclass[11pt,a4paper]{article}
\usepackage{fontspec}
\setmainfont{Times New Roman}
\usepackage[russian]{babel}
\usepackage[margin=22mm]{geometry}
\begin{document}
\section{Профиль и протокол}
Профиль P3adf02. Глубины 8, 24, 48, 96; ширина 256; гауссовская инициализация с $\sigma=0.02$; SGD с импульсом 0.9. CIFAR-10 разрешён методическими указаниями. Общий механизм экспериментов сохранён из ЛР0.
\section{E1: гипотеза до эксперимента}
По формулам (1.4)--(1.8) ожидаются множители 1 для He, 0.5 для Xavier и 0.0512 для гауссовской схемы. Корреляции конечной сети могут нарушать независимость весов и сигналов.
\section{E2: гипотеза до эксперимента}
По формуле (1.13) единичная ветвь даёт экспоненциальный рост, а масштаб $1/\sqrt{L}$ ограничивает накопление среднего квадрата. Точность и обучаемость будут измерены раздельно.
\end{document}
